\paragraph{Uniform finite first-order certificate.}
For the scaled-fiber parameters $N=9p^2$, $K=3$, and $A=4p$,
there is also an explicit finite certificate, independent of the limiting
rate curve. In the derivative-weighted support of~\cite[Eq.~(60)]{dkt2026}, take
\[
 m=6,\qquad B_\partial=3,\qquad
 B=B_{\rm jet}=12p,\qquad H=B_Z=40B.
\]
Thus the support consists of $X^xY_0^aY_1^b$ with
$b\le3$, $a+b\le B$, and $x+2a+b<24p$.
The exact coefficient count in~\cite[Eq.~(61)]{dkt2026} is
\[
 G=\sum_{b=0}^{3}\sum_{t=b}^{B}(2B-2t+b)
   =\sum_{b=0}^{3}B(B-b+1)
   =4B^2-2B=576p^2-24p.
\]
The six local-rank contributions in~\cite[Eq.~(62)]{dkt2026}
are $4,8,12,15,14,9$, so $R=62$.
Every column has remaining challenge budget at least $H-B$,
and every local row has budget at most $H$. Consequently the graded
challenge test of~\cite[Proposition~5.10]{dkt2026} follows from
\begin{align*}
 (H-B+1)G-(H+1)NR
 &=B(144p^2-936p)+18p^2-24p\\
 &>0 \qquad (p\ge7).
\end{align*}
In particular $G>NR$. This constructs a sound first-order interpolant,
nonzero at every challenge specialization, uniformly over all evaluation
sets of this size and all affine received lines. Its actual derivative
variable degree is at most three, so the reconstruction and counting
hypothesis is only
\[
 \operatorname{char}(\mathbb F)=p>\max\{K-1,B_\partial\}=3.
\]
Thus the stated family satisfies an explicit finite first-order test;
no characteristic bound depending on the growing total jet cap is needed.
The corresponding uniform list and full-agreement-set MCA bounds follow
from the counting transfers of~\cite[Lemma~5.6]{dkt2026}; we do not optimize
their numerical constants here.
